\TrangKhongSo{KẾT LUẬN}

\section*{Các kết quả đạt được}

Đồ án đã xây dựng hệ thống FlexSDS phục vụ mô hình ngôn ngữ lớn tương thích OpenAI, tích hợp giải mã suy đoán EAGLE với bộ điều khiển ngân sách nháp thích ứng, và đánh giá hệ thống trên quy trình xử lý hoá đơn tiếng Việt. Đối chiếu với các mục tiêu đề ra, kết quả đạt được như sau.

Về mục tiêu xây dựng các làn đối chứng và bộ đo, nhóm đã xây dựng năm làn phục vụ trên vLLM 0.30.0 (giải mã thường, mô hình nháp 0,5B với ngân sách cố định và thích ứng, EAGLE-3.1 với ngân sách cố định và thích ứng) cùng hợp đồng thực thi tương đồng để so sánh công bằng, và một bộ đo tự động với workload đóng băng, kiểm tra môi trường, khởi động khớp tải, lặp lại và tổng hợp có truy vết. Mục tiêu này hoàn thành.

Về mục tiêu bộ điều khiển, nhóm đã hiện thực bộ điều khiển ba lớp (bảo vệ bộ nhớ, chọn theo chi phí đo, quy tắc ngưỡng dự phòng) và tích hợp vào bộ lập lịch bất đồng bộ của vLLM mà không sửa tiêu chí chấp nhận, kèm cơ chế quy phản hồi đúng nguồn có kiểm tra mã nguồn của engine. Do giới hạn của vLLM 0.30.0, việc điều khiển được thực hiện trên số token nháp vô hướng thay vì cấu trúc cây như dự kiến ban đầu. Về mặt chức năng, mục tiêu này hoàn thành; về mặt hiệu quả, bộ điều khiển chưa đạt kỳ vọng như trình bày dưới đây.

Về mục tiêu đóng gói, công cụ dòng lệnh \texttt{flexsds} đã cung cấp đầy đủ luồng nhận diện phần cứng, kiểm tra tương thích, đo hồ sơ chi phí, khởi chạy cục bộ hoặc trên Modal và xác minh sẵn sàng bằng một lần sinh văn bản thực. Nhóm cũng đã huấn luyện và công bố đầu nháp EAGLE-3.1 cho mô hình đích, và xây dựng web console để so sánh trực quan các endpoint trên quy trình hoá đơn. Mục tiêu này hoàn thành.

Về mục tiêu đánh giá, nhóm đã thực hiện bảy chiến dịch trên ba loại GPU. Câu trả lời cho các câu hỏi nghiên cứu như sau. Với RQ1, hai phiên bản bộ điều khiển thích ứng chưa vượt được cấu hình cố định: phiên bản dùng quy tắc ngưỡng kém mô hình nháp cố định ở hầu hết các ô, và phiên bản dùng hồ sơ chi phí kém EAGLE với ngân sách cố định ở mọi ô trên A10G. Với RQ2, tăng tốc âm chưa được loại bỏ vì việc chọn không sinh nháp vẫn để bộ nháp chạy, nên chi phí một bước không trở về mức của giải mã thường; không có lỗi tràn bộ nhớ nào xảy ra, nhưng lớp bảo vệ bộ nhớ chưa được kiểm chứng dưới áp lực thực. Với RQ3, giải mã suy đoán có lợi rõ rệt ở tải thấp và vừa; đầu nháp EAGLE-3.1 với ngân sách cố định nhỏ tăng tốc $1{,}3$--$1{,}7\times$ trên cả ba loại tác vụ khi số yêu cầu đồng thời không quá 8, trong khi ở tải cao không cấu hình suy đoán nào vượt được giải mã thường. Về tính lossless, đầu ra của các làn suy đoán khớp từng token với giải mã thường ở một phần các prompt; mọi điểm lệch được truy nguyên tới các vị trí hoà điểm trong giới hạn độ chính xác bfloat16, và chất lượng trích xuất nghiệp vụ không thay đổi.

\section*{Đóng góp của đồ án}

Đóng góp chính của đồ án gồm hệ thống FlexSDS có thể cài đặt và sử dụng; đầu nháp EAGLE-3.1 cho Qwen2.5-7B-Instruct; bộ số liệu thực nghiệm có kiểm soát về giải mã suy đoán trên GPU phổ thông theo thang mức đồng thời; và phân tích định lượng, dựa trên telemetry từng bước, về các điều kiện mà một bộ điều khiển thích ứng cần thoả mãn. Cụ thể, đồ án chỉ ra ba điều kiện: hồ sơ chi phí phải đo thời gian đầu--cuối của một bước phục vụ, bao gồm chi phí cố định của bộ lập lịch, phần lấy mẫu và cơ chế điều khiển; lựa chọn ``tắt suy đoán'' phải thực sự loại bỏ chi phí của bộ nháp; và mô hình tỉ lệ chấp nhận phải được hiệu chỉnh theo vị trí. Các điều kiện này không bộc lộ khi chỉ đánh giá ở một yêu cầu trên GPU cao cấp.

\section*{Hạn chế}

Đồ án còn các hạn chế sau. Bộ điều khiển thích ứng chưa đạt mục tiêu tốc độ và chưa loại bỏ được tăng tốc âm. Hồ sơ chi phí hợp lệ mới chỉ có cho GPU A10G, nên bộ điều khiển theo chi phí chưa được đánh giá trên T4 và L4. Các chiến dịch được chạy trên ba loại GPU khác nhau do khả năng cấp phát của nền tảng, nên không có một GPU tham chiếu chung cho mọi so sánh; các chiến dịch trước R6 còn lẫn hiệu ứng của chế độ thực thi giữa làn cố định và làn thích ứng. Thực nghiệm chỉ gồm một mô hình đích, workload có quy mô nhỏ và thuộc một miền nghiệp vụ. Đầu ra không khớp từng token với giải mã thường trên phần cứng thực, và phép kiểm tra với độ chính xác số học cao hơn chưa được thực hiện. Lớp bảo vệ bộ nhớ chưa được kiểm chứng trong điều kiện áp lực bộ nhớ thực. Cuối cùng, chiến dịch R6 được thực hiện sau ngày đóng băng mã nguồn và kết quả của nó đang chờ nhóm xác nhận chính thức.

\section*{Hướng phát triển}

Từ các phân tích ở Chương~\ref{chap:thu-nghiem}, nhóm đề xuất các hướng phát triển sau. Thứ nhất, đo lại hồ sơ chi phí theo thời gian đầu--cuối của một bước phục vụ trên đúng đường thực thi của chế độ thích ứng, để mô hình chi phí bao gồm các chi phí cố định. Thứ hai, hiệu chỉnh ước lượng tỉ lệ chấp nhận theo vị trí trước khi đưa vào công thức kỳ vọng số token. Thứ ba, bổ sung cơ chế tắt suy đoán ở cấp máy chủ, chẳng hạn định tuyến yêu cầu giữa một máy chủ giải mã thường và một máy chủ có bộ nháp theo mức tải, thay vì chỉ đặt số token nháp bằng 0. Thứ tư, khi engine hỗ trợ truyền cấu trúc cây nháp tuỳ ý, mở rộng bộ điều khiển sang chọn cấu trúc cây theo chi phí. Thứ năm, kiểm tra tính lossless với chế độ tính bất biến theo batch hoặc độ chính xác số học cao hơn. Cuối cùng, mở rộng thực nghiệm sang các GPU, mô hình đích và workload khác, đo hồ sơ cho T4 và L4, và thử nghiệm lớp bảo vệ bộ nhớ dưới tải ngữ cảnh dài.
